\documentclass[11pt]{article}
\usepackage[margin=1in]{geometry}
\usepackage[T1]{fontenc}
\usepackage[utf8]{inputenc}
\usepackage{lmodern}
\usepackage{xcolor}
\usepackage{hyperref}
\usepackage{enumitem}
\usepackage{parskip}

\definecolor{slashblue}{RGB}{48,92,222}
\definecolor{slashdark}{RGB}{20,28,45}

\hypersetup{
    colorlinks=true,
    linkcolor=slashblue,
    urlcolor=slashblue
}

\title{\textbf{AI Firewall Proxy Repository Overview}}
\author{}
\date{}

\begin{document}

\maketitle
\thispagestyle{empty}

\section*{Overview}

This repository implements an \textbf{AI firewall proxy} for agentic applications. Its purpose is to sit between an AI agent and the underlying model provider, inspect incoming requests, and decide whether they should be allowed, warned on, denied, or bypassed according to configurable security and scope policies.

The project is designed for teams building AI agents that communicate with OpenAI-compatible APIs or other HTTP-based workflows. Instead of changing the core agent logic, developers point their agent to this proxy. The proxy then becomes the security control plane for traffic inspection, forwarding, blocking, logging, and operational visibility.

\section*{Purpose}

The main goal of the repository is to provide a reusable and developer-friendly security layer for AI systems. In practical terms, it helps:

\begin{itemize}[leftmargin=1.2em]
    \item prevent out-of-scope requests from being sent to an agent,
    \item detect prompt injection attempts before they reach the model,
    \item preserve the normal application flow when requests are allowed,
    \item centralize logging, review, and policy tuning,
    \item support live operational control through a dashboard and runtime toggles.
\end{itemize}

This is especially useful for domain-specific agents, such as legal, travel, or customer-support assistants, where the system should answer only within an intended domain and reject unsafe or manipulated prompts.

\section*{What The Repository Contains}

The repository contains four main functional parts:

\begin{enumerate}[leftmargin=1.4em]
    \item \textbf{FastAPI proxy layer} \\
    Implemented primarily in \texttt{main.py}, this is the request gateway. It receives incoming API calls, applies firewall logic to protected routes, forwards allowed requests to the configured upstream model backend, and returns the resulting model response to the calling agent.

    \item \textbf{Firewall logic and configuration} \\
    Implemented in \texttt{firewall\_proxy/firewall.py} and \texttt{firewall\_proxy/config.py}. This layer performs the security checks and loads agent-specific policy definitions from \texttt{config.yaml}. Each agent can have its own description, allowed examples, and denied examples, making the system suitable for multi-agent or multi-tenant use cases.

    \item \textbf{Logging and observability} \\
    Implemented in \texttt{firewall\_proxy/log\_store.py} and the Streamlit dashboard in \texttt{dashboard.py}. Requests and decisions are logged to SQLite, and the dashboard provides a visual interface for inspecting decisions, scores, recent events, and live protection state.

    \item \textbf{Runtime control and live toggles} \\
    Implemented in \texttt{firewall\_proxy/runtime\_state.py} and exposed through the dashboard. This allows operators to enable or bypass firewall enforcement globally without changing the agent configuration or restarting the whole stack.
\end{enumerate}

\section*{Current Detection Approach}

At present, the repository uses two primary checks:

\begin{itemize}[leftmargin=1.2em]
    \item a \textbf{scope check} based on sentence embeddings using \texttt{sentence-transformers/all-MiniLM-L6-v2}, and
    \item a \textbf{prompt injection classifier} loaded locally from the repository for direct prompt injection detection.
\end{itemize}

The scope check helps ensure that requests match the purpose of the agent. For example, a Massachusetts housing-law assistant should answer landlord-tenant questions, but not unrelated travel or cybersecurity prompts. The prompt injection classifier is used to detect malicious attempts such as instruction override, hidden prompt extraction, and jailbreak-style inputs.

\section*{Planned Enhancement: Llama Guard Prompt 22M}

A near-term enhancement to this repository is the integration of \textbf{Llama Guard Prompt 22M} for \textbf{indirect prompt injection detection}. This addition is intended to strengthen the firewall against attacks that do not only target the visible user prompt, but instead try to manipulate tool use, trigger unsafe actions, or escalate privileges through retrieved or external content.

In presentation terms, the planned Llama Guard Prompt 22M integration is meant to cover risks such as:

\begin{itemize}[leftmargin=1.2em]
    \item \textbf{tool misuse},
    \item \textbf{privilege escalation},
    \item \textbf{indirect prompt injection embedded in documents or retrieved context},
    \item \textbf{unsafe action triggering through external instructions}.
\end{itemize}

This capability is \textbf{not yet fully integrated in the current repository state}, but it is the next planned security component and aligns directly with the project roadmap.

\section*{Overall Value}

In summary, this repository provides a practical foundation for deploying AI agents behind a reusable security layer. It combines request filtering, prompt-injection defense, scope validation, logging, dashboard-based monitoring, and live enforcement controls in a single project. Its broader purpose is to make agent deployments safer, more observable, and easier to govern without requiring developers to redesign their application workflows.

\vfill
\begin{center}
\textcolor{slashdark}{\small Prepared as a short technical overview for presentation purposes.}
\end{center}

\end{document}
